\documentclass[11pt]{article}

% ======================
% Engine: LuaLaTeX
% ======================
\usepackage[margin=1in]{geometry}
\usepackage{fontspec}
\setmainfont{Latin Modern Roman}
\usepackage{microtype}

\usepackage{amsmath,amssymb,amsthm,mathtools}
\usepackage{tikz-cd}
\usepackage{hyperref}
\usepackage[nameinlink]{cleveref}
\usepackage{listings}

\hypersetup{
colorlinks=true,
linkcolor=blue,
citecolor=blue,
urlcolor=blue
}

\lstset{
basicstyle=\ttfamily\small,
columns=fullflexible,
frame=single,
xleftmargin=2pt,
xrightmargin=2pt,
breaklines=true
}

\newtheorem{theorem}{Theorem}
\newtheorem{proposition}{Proposition}
\newtheorem{lemma}{Lemma}
\newtheorem{definition}{Definition}
\newtheorem{remark}{Remark}

\title{Algorithmic Structure Towers: Unifying Bourbaki's Structuralism and Computable Martingale Theory in Lean 4}

\author{
su\
(TeX generated by ChatGPT; Lean code generated by OpenAI Codex)%
}
\date{\today}

\begin{document}
\maketitle

% ==========================================================
\begin{abstract}
We present \emph{algorithmic structure towers}, a Lean~4 formal framework that reconstructs Bourbaki-style structuralism as a stratified object equipped with a canonical ``first appearance'' operator \texttt{minLayer}. % StructureTowerWithMin definition source. 
The key design goal is to make \emph{abstraction computational}: by restricting probabilistic development to finite sample spaces and rational-valued processes, definitional semantics and executable evaluation (\texttt{#eval}) coincide.

Our main conceptual contribution is a uniform mapping from probabilistic time concepts to tower-theoretic primitives: \emph{stopping times} are recovered via the minimal layer at which stopping sets become measurable, using a filtration-as-tower bridge and a \texttt{Nat.find}-based minimality principle. % Filtration minLayer + stopping-time bridge.  
On the formal side, we outline a five-layer architecture culminating in a constructive Optional Stopping Theorem (OST) proof pattern based on telescoping decompositions, and we provide executable sanity checks and derived constructions (e.g.\ stopped processes and bounded stopping). % stopped martingale construction. 
\end{abstract}

% ==========================================================
\section{Introduction}

Bourbaki's structural program emphasized \emph{species of structures} and their transport across domains. The modern proof-assistant ecosystem offers an opportunity to re-express this program as a reusable, compositional, and machine-checked infrastructure, where the \emph{shape} of an argument becomes a programmable artifact.

A persistent obstacle in formal probability is the mismatch between mathematical definitions (often measure-theoretic, infinitary, and real-valued) and computational meaning. Our approach makes the mismatch explicit and then removes it by design: we work in a discrete, finitistic setting where measurability and expectation are \emph{decidable or computably reducible} (finite carriers; rational arithmetic). The result is a ``living theory'' where definitions are not merely symbolic but executable.

This paper contributes:

\begin{itemize}
\item A tower-based reconstruction of Bourbaki structuralism in Lean~4 with a canonical minimality operator \texttt{minLayer}. % Bourbaki/structure tower hierarchy and minLayer role. 
\item A categorical stratification of morphisms capturing varying levels of strictness (existential layer preservation vs.\ explicit index maps vs.\ \texttt{minLayer}-preservation), enabling a precise match to probabilistic notions. % Cat_D/Cat_le hierarchy motivation.  
\item A filtration-as-tower bridge with a \texttt{Nat.find}-based \texttt{minLayer} for events, and a stopping-time development that isolates the stopping filtration core and stopped processes. % sigma-algebra minLayer and stopping-time module goals.  
\item A constructive OST proof blueprint based on telescoping decompositions and bounded stopping, together with executable checks.
\end{itemize}

% ==========================================================
\section{Theory: The Structure Tower}

\subsection{The core object with minimal layer}

We begin with the minimal structure needed to express stratified ``stages'' and a canonical first stage of membership.

\begin{definition}[Structure tower with minimal layer]
A \emph{structure tower with minimal layer} consists of:
[
(X,\ I,\ \le,\ L: I\to \mathcal P(X),\ \mathrm{covering},\ \mathrm{monotone},\ \minLayer: X\to I)
]
such that every element lies in some layer (covering), layers are monotone in the index order, and \texttt{minLayer} picks a minimal index witnessing membership.
\end{definition}

In Lean~4 (schematically), this is:

\begin{lstlisting}[language={},caption={Lean skeleton: \texttt{StructureTowerWithMin}}]
structure StructureTowerWithMin where
carrier : Type u
Index : Type v
[indexPreorder : Preorder Index]
layer : Index → Set carrier
covering : ∀ x : carrier, ∃ i : Index, x ∈ layer i
monotone : ∀ {i j : Index}, i ≤ j → layer i ⊆ layer j
minLayer : carrier → Index
minLayer_mem : ∀ x, x ∈ layer (minLayer x)
minLayer_minimal : ∀ x i, x ∈ layer i → minLayer x ≤ i
\end{lstlisting}
% Exact field list appears in the CAT2 StructureTower explanation. 

A representative example is the integer tower stratified by absolute value, where \texttt{minLayer k = |k|} and membership is decidable, allowing evaluation-level sanity checks. % integer tower and #eval lines. 

\subsection{Rank normal form: towers as canonical ranks}

A central simplification is that, over $\mathbb N$-indexed towers, the tower can be put into a \emph{rank normal form}. Intuitively, \texttt{minLayer} \emph{is} the rank, and the $n$-th layer is the set of points of rank $\le n$.

\begin{proposition}[Rank-to-tower constructor]
Given a rank function $\rho : X \to \mathrm{WithTop}(\mathbb N)$ satisfying a mild cover condition, one can construct a tower with
[
L(n)={x\mid \rho(x)\le n},\qquad \minLayer(x)=\mathrm{Nat.find}(\mathrm{cover}(x)).
]
\end{proposition}

This construction is implemented as \texttt{towerFromRank}. % towerFromRank definition and minLayer via Nat.find. 

\begin{remark}
This rank normal form is not merely a convenience: it is a Bourbaki-compatible \emph{canonicalization} principle. It explains why repeated proof patterns become mechanically uniform once layers are expressed as ${x\mid \rho(x)\le n}$. % proof-pattern claim and L(n)={x | ρ(x) ≤ n}. 
\end{remark}

\subsection{Morphisms: strict vs.\ upper-bound preservation}

A structural program needs a principled notion of structure-preserving map. In the tower setting, there are multiple useful strengths.

\paragraph{(i) Existential layer preservation: \texttt{Cat_D}.}
The weakest and most flexible notion only requires that the image of a layer lands in \emph{some} layer.

\begin{lstlisting}[language={},caption={Lean skeleton: \texttt{Cat_D} morphism}]
structure HomD (T T' : TowerD) where
map : T.carrier → T'.carrier
map_layer : ∀ i : T.Index, ∃ j : T'.Index,
Set.image map (T.layer i) ⊆ T'.layer j
\end{lstlisting}
% HomD definition and motivation for probability applications.  

This is appropriate for probabilistic transport, where one often only needs that a transformed event becomes measurable at \emph{some} time (not necessarily a uniquely determined one). % motivation: measurable at some σ-algebra. 

\paragraph{(ii) Explicit index maps: \texttt{Cat_le}.}
A refinement carries an explicit monotone index map $\varphi$ (time change), making order preservation visible.

\begin{lstlisting}[language={},caption={Lean skeleton: \texttt{Cat_le} morphism}]
structure HomLe (T T' : TowerD) where
map : T.carrier → T'.carrier
indexMap : T.Index → T'.Index
indexMap_mono : Monotone indexMap
layer_preserving : ∀ i x, x ∈ T.layer i → map x ∈ T'.layer (indexMap i)
\end{lstlisting}
% Cat_le: explicit indexMap and monotonicity.  

\paragraph{(iii) \texttt{minLayer}-preserving morphisms: \texttt{Cat2}.}
A still stricter level additionally requires preservation of the canonical minimal layer, enabling uniqueness of induced index behavior and supporting universal constructions. % Cat2 role described alongside Cat_D/Cat_le.  

\begin{remark}[Relevance to martingales]
The distinction between equality-preserving and inequality-preserving laws in probability (martingale vs.\ submartingale) suggests an analogy with strict vs.\ upper-bound morphisms, but our current development uses this mainly as an \emph{architectural} separation: strictness is crucial for universal properties, while probabilistic transport benefits from existential flexibility. % Cat_D is ``most flexible'' for probability; hierarchy. 
\end{remark}

% ==========================================================
\section{Application: Computable Martingales}

\subsection{A five-layer discrete architecture}

Our probabilistic formalization is designed as a tower-instantiation pipeline:

\begin{enumerate}
\item \textbf{Decidable events} on finite sample spaces (membership decidable).
\item \textbf{Decidable finite algebras} of events (finite $\sigma$-algebra analog).
\item \textbf{Decidable filtrations} as monotone families of algebras.
\item \textbf{Computable stopping times} built from computable combinators (e.g.\ \texttt{const}, \texttt{min}, \texttt{max}), with boundedness and derived constructions.
\item \textbf{Algorithmic martingales and OST}, using rational-valued processes with exact expectation computation and executable checks.
\end{enumerate}

The interface point where the general tower theory directly constrains probability is Layer~3: a filtration is precisely a tower of ``observable events at time $n$'', and \texttt{minLayer} becomes the earliest observation time of an event.

\subsection{Filtrations as towers and \texttt{minLayer} for events}

In the measure-theoretic (but still discrete-indexed) side of the development, a filtration core is a monotone family of measurable spaces, and \texttt{minLayer} for an event $A$ is defined by \texttt{Nat.find}. Concretely:

\begin{lstlisting}[language={},caption={Lean: \texttt{minLayer} for $\mathbb N$-filtrations}]
noncomputable
def minLayer (F : NatFiltration Ω) (A : Set Ω)
(hA : ∃ n : ℕ, @MeasurableSet Ω (F.𝓕 n) A) : ℕ :=
Nat.find hA
\end{lstlisting}
% minLayer for filtrations and its two key lemmas. 

This comes with:
(i) measurability at the returned layer, and (ii) minimality among all measurable layers. % minLayer_measurable and minLayer_le_of_measurable. 

The stopping-time module then explicitly promotes a filtration to a tower and positions \texttt{minLayer} as the primitive that measures ``when an event becomes observable.'' % towerOf bridge. 

\subsection{Stopping times via stopping sets and minimal observability}

A classical stopping time $\tau$ (discrete time) is specified by measurability of stopping sets ${\tau \le n}$ at time $n$:

\begin{lstlisting}[language={},caption={Lean: classical stopping time}]
structure StoppingTime (ℱ : Filtration Ω) where
τ : Ω → ℕ
measurable : ∀ n, @MeasurableSet Ω (ℱ.base.𝓕 n) {ω | τ ω ≤ n}
\end{lstlisting}
% StoppingTime definition in StoppingTime_MinLayer. 

The same file develops the \emph{stopped filtration core} via truncated stopping times and proves that each stopped layer embeds in the original filtration, as well as measurability of stopping sets in the stopped core. % stoppedFiltrationCore, monotonicity, and stoppingSet measurable lemma.  

\paragraph{Key insight (tower-theoretic re-interpretation).}
Given an event-former $A_n := {\tau \le n}$, the tower perspective treats the mapping
[
n \longmapsto A_n
]
as a stratified object whose minimality statements can be expressed as inequalities involving \texttt{minLayer}. Informally, ``$\tau$ is the first time when the stopping set becomes measurable'' is encoded by the minimality lemmas associated to \texttt{Nat.find}. % minLayer minimality lemma for filtrations. 

\subsection{Stopped processes as pure functions}

To support optional stopping proofs, we define the stopped process purely at function level:
[
X^{\tau}*n(\omega)=X*{\min(n,\tau(\omega))}(\omega).
]

\begin{lstlisting}[language={},caption={Lean: stopped process and basic laws}]
def stopped {β : Type*} (X : ℕ → Ω → β) (τ : Ω → ℕ) : ℕ → Ω → β :=
fun n ω => X (Nat.min n (τ ω)) ω

lemma stopped_eq_of_le ... : stopped X τ n ω = X n ω := ...
lemma stopped_const_of_ge ... : stopped X τ m ω = stopped X τ n ω := ...
\end{lstlisting}
% stopped definition and equalities. 

A thin wrapper \texttt{StoppedProcess} packages a filtration, an underlying process, and a stopping time, providing convenient lemmas such as `before stopping equals original'' and `after stopping is constant.'' % StoppedProcess wrapper and lemmas. 

% ==========================================================
\section{Main Result: Optional Stopping Theorem}

\subsection{Statement (bounded discrete-time OST)}

We focus on the bounded case (the setting most compatible with executability and finite models).

\begin{theorem}[Bounded Optional Stopping (discrete time, blueprint)]
Let $(M_n)*{n\in\mathbb N}$ be a martingale with respect to a filtration $(\mathcal F_n)$ on a finite sample space, and let $\tau$ be a stopping time bounded by $K$ (i.e.\ $\tau(\omega)\le K$ for all $\omega$). Then
[
\mathbb E[M*{\tau}] = \mathbb E[M_0].
]
Moreover, the stopped process $(M^{\tau}_n)$ is itself a martingale.
\end{theorem}

The Lean development contains the key structural step: from a martingale $M$ and a bounded stopping time satisfying measurability of stopping sets, we construct a new martingale whose process is the stopped process. % stoppedProcess_martingale_of_bdd. 

\subsection{Proof sketch via telescoping decomposition}

The proof follows the standard discrete telescoping pattern, phrased in a way that aligns with mechanization.

Define increments $D_n := M_{n+1}-M_n$. Then for any $N$,
[
M_{\tau\wedge N} - M_0 ;=; \sum_{n=0}^{N-1} D_n \cdot \mathbf 1_{{\tau>n}}.
]
Taking expectations, it suffices to show each summand has expectation $0$:
[
\mathbb E!\left[D_n \cdot \mathbf 1_{{\tau>n}}\right]=0.
]
Because ${\tau>n}\in\mathcal F_n$ (equivalently, ${\tau\le n}\in\mathcal F_n$), and the martingale condition implies $\mathbb E[D_n\mid \mathcal F_n]=0$, we obtain
[
\mathbb E[D_n \cdot \mathbf 1_{{\tau>n}}]
=\mathbb E!\left(\mathbb E[D_n\mid \mathcal F_n]\cdot \mathbf 1_{{\tau>n}}\right)=0.
]
Boundedness lets us take $N\ge K$ and conclude $\tau\wedge N=\tau$.

\paragraph{Where towers enter.}
The only genuinely ``time-sensitive'' measurability premise is ${\tau\le n}\in\mathcal F_n$; the tower formulation highlights that this is precisely about the layer at which a stopping set appears, making \texttt{minLayer} a natural organizational primitive for stopping-time reasoning. % stopping-time measurable sets + filtration minLayer idea.  

\subsection{Executable validation}

A core motivation is that the formal objects admit evaluation-level sanity checks. For instance, the integer absolute-value tower computes minimal layers exactly:

\begin{lstlisting}[language={},caption={Lean: evaluation-level checks (example)}]
#eval intAbsTower.minLayer (0 : ℤ)      -- 0
#eval intAbsTower.minLayer (5 : ℤ)      -- 5
#eval intAbsTower.minLayer (-3 : ℤ)     -- 3
\end{lstlisting}
% These #eval lines appear in DecidableStructureTower_Examples. 

In the probability pipeline, the analogous goal is that expectations of rational-valued processes over finite carriers reduce to finite sums and therefore become computable objects (and testable by evaluation). The bounded optional stopping construction then becomes not only provable but \emph{checkable} on concrete finite models.

% ==========================================================
\section{Conclusion}

We introduced algorithmic structure towers as a unifying Lean~4 framework that (i) reconstructs Bourbaki-style structuralism as stratified towers with a canonical \texttt{minLayer}, (ii) equips the theory with a categorical stratification of morphisms matching application needs, and (iii) applies the resulting infrastructure to a computable martingale program culminating in optional stopping reasoning.

The conceptual payoff is the identification of stopping-time measurability as a tower-theoretic minimality phenomenon: ``time'' becomes \emph{the first layer of observability}. The engineering payoff is a development style where abstract definitions, mechanized proofs, and executable computations reinforce each other.

\paragraph{Outlook.}
Immediate extensions include strengthening categorical results (limits/colimits, adjunctions), expanding applications (Cantor--Bendixson style ranks, closure-comonad viewpoints), and analyzing the computational complexity of \texttt{minLayer}-driven constructions. % future directions and CB-rank link to minLayer.  

% ==========================================================
\section*{Acknowledgments}
This manuscript's \LaTeX{} source was generated using ChatGPT, and the associated Lean~4 codebase was generated using OpenAI Codex under human supervision, then iteratively refined to satisfy compilation and proof obligations.

% ==========================================================
\begin{thebibliography}{9}

\bibitem{BourbakiSets}
N.\ Bourbaki.
\newblock \emph{'{E}l'{e}ments de math'{e}matique: Th'{e}orie des ensembles}.
\newblock Hermann.

\bibitem{MacLane}
S.\ Mac~Lane.
\newblock \emph{Categories for the Working Mathematician}.
\newblock Springer, 1971.

\bibitem{Awodey}
S.\ Awodey.
\newblock \emph{Category Theory}.
\newblock Oxford University Press.

\bibitem{Mathlib}
The mathlib Community.
\newblock \emph{The Lean Mathematical Library (mathlib)}.

\end{thebibliography}

\end{document}
